% ECONOMICS_PROBLEM: {"id": "I14-603", "title": "Socioeconomic health gradients", "jel_code": "I14", "source_id": "OP-0056"}
% !TeX program = xelatex
\documentclass[11pt,a4paper]{article}
\usepackage[margin=23mm]{geometry}
\usepackage{fontspec}
\setmainfont{Latin Modern Roman}
\usepackage{amsmath,amssymb,mathtools}
\usepackage{xcolor,enumitem,booktabs,longtable,array,xurl,hyperref}
\definecolor{muted}{HTML}{53636C}
\hypersetup{colorlinks=true,urlcolor=blue,linkcolor=blue}
\setlength{\parindent}{0pt}
\setlength{\parskip}{5pt}
\newcommand{\R}{\mathbb R}
\newcommand{\E}{\mathbb E}
\newcommand{\N}{\mathbb N}
\newcommand{\ind}{\mathbf 1}
\newcommand{\Var}{\operatorname{Var}}
\newcommand{\Cov}{\operatorname{Cov}}
\newcommand{\supp}{\operatorname{supp}}
\DeclareMathOperator*{\argmin}{arg\,min}
\DeclareMathOperator*{\argmax}{arg\,max}
\newcommand{\entrymeta}[3]{{\sffamily\small\color{muted}\textbf{\texttt{#1}}\quad|\quad #2\par #3\par}}
\newcommand{\Problemref}[1]{\texttt{#1}}
\newcommand{\JELref}[2]{\texttt{#2} (JEL \texttt{#1})}
\begin{document}
\section*{Socioeconomic health gradients}
\textbf{Problem ID:} \texttt{I14-603}\quad \textbf{JEL:} \texttt{I14}\par
% BEGIN ECONOMICS STATEMENT
\label{problem:OP-0056}
\label{atlas:prob:H6}
\noindent\textbf{Original atlas identifier:} H6; entry 56. \textbf{Field:} Health economics. \textbf{Original tractability label:} Medium.

\noindent\textbf{Classification:} Parameterized theoretical/empirical research agenda; not a certified unsolved theorem.

\subsection*{Definitions and assumptions}
Fix a birth cohort and dates $t=0,\ldots,T$. Let $Y_t$ be disposable income, $H_t$ a specified health outcome, $X_t$ observed education/environment/employment histories, and $U_t$ latent determinants. A causal model $m$ gives ordered structural transition equations for these variables with a declared within-period timing convention; feedback across periods is allowed. An income intervention $a$ specifies a transfer path and its financing, not merely conditioning on high observed income. A health intervention $b$ specifies a clinical or disability-prevention change. Define counterfactual paths $H_T(a)$ and $Y_T(b)$ with consistency, positivity, and interference conditions explicitly stated for the target population. Baseline income--health covariance is descriptive; it is not assumed to possess a unique additive causal decomposition.

\subsection*{Formal research question}
\emph{Editorial formulation: the mathematical target below is not a theorem quoted from the references.}

Identify intervention-specific causal targets
\[
 \Delta_H(a,a_0)=\mathbb E[H_T(a)-H_T(a_0)],\qquad
 \Delta_Y(b,b_0)=\mathbb E[Y_T(b)-Y_T(b_0)].
\]
Here $a_0,b_0$ are defined baseline interventions. Determine the assumptions under which randomized income variation, lotteries, sibling designs, or longitudinal instruments identify these targets, and bound them when exclusion or sequential-exchangeability conditions fail. If the research goal retains a “fraction” of the observed gradient, define an explicit attribution functional $A_m$ (including order of pathway removal and treatment of interactions) and study its identified set. Do not equate a covariance share with either causal target or assert that forward and reverse shares sum to one.

\subsection*{Interpretation and scope}
Income can affect nutrition, housing, stress, and care, while health affects productivity, labor supply, and family resources. A shock to income may change employment conditions or insurance simultaneously; its policy effect is meaningful, but interpreting it as a pure income channel requires further exclusions. Conditioning on education or occupation may block mediating pathways or induce collider bias, depending on temporal ordering. Cohort survival and selective measurement also change who contributes to a measured gradient. Sibling comparisons remove only specified shared factors and can amplify nonshared measurement error. Descriptive mortality trends motivate mechanisms but do not identify an income intervention. The rigorous extension concerns well-defined dynamic interventions and pathway attribution under explicit models rather than a universal numerical fraction of correlation.

\subsection*{Source audit and status}
[S1] is health-capital theory, [S2] a mechanisms review, and [S3] descriptive morbidity/mortality analysis. None identifies the requested two-direction “fraction” without additional causal definitions. The original wording is therefore corrected as a partly ill-posed estimand, not advertised as an unresolved algebraic theorem.

\subsection*{References}
\begin{enumerate}[label={[S\arabic*]},leftmargin=*,itemsep=0.6em]
\item Michael Grossman (1972). On the Concept of Health Capital and the Demand for Health. Journal of Political Economy 80(2), 223--255. DOI: 10.1086/259880.
\url{https://www.journals.uchicago.edu/doi/10.1086/259880}
\newline\emph{Locator and support limits:} Article: health-capital theory; no empirical decomposition of income-health covariance.
\item David M. Cutler, Adriana Lleras-Muney, and Tom Vogl (2011). Socioeconomic Status and Health: Dimensions and Mechanisms. In Sherry Glied and Peter C. Smith (eds.), The Oxford Handbook of Health Economics, chapter 7. Oxford University Press.
\url{https://academic.oup.com/edited-volume/28339/chapter-abstract/215132856}
\newline\emph{Locator and support limits:} Chapter: socioeconomic gradients and mechanisms. NBER WP14333 (2008) is the earlier working-paper version.
\item Anne Case and Angus Deaton (2015). Rising Morbidity and Mortality in Midlife among White Non-Hispanic Americans in the 21st Century. Proceedings of the National Academy of Sciences 112(49), 15078--15083. DOI: 10.1073/pnas.1518393112.
\url{https://accase.scholar.princeton.edu/publications/rising-morbidity-and-mortality-midlife-among-white-non-hispanic-americans-21st}
\newline\emph{Locator and support limits:} Author publication record and paper: descriptive morbidity/mortality trends; not an income intervention or causal covariance decomposition.
\end{enumerate}

\subsection*{Common conventions: Source mathematical conventions}

Unless an entry states otherwise, agent and firm sets are finite. State and action spaces are measurable, strategies and outcome maps are measurable, probability laws are specified, and expectations exist for the targets under discussion. Infinite-horizon discount factors lie in $(0,1)$ and discounted objectives are finite. Norms, tolerances and complexity input sizes must be specified for computational comparisons. Constants or primitives that appear locally are inputs, not empirical facts asserted by this catalogue.

\textbf{Identification.} A statistical model $m\in\mathcal M$ implies an observable law $P_m$ and target $\theta(m)$. At law $P$, its identified set is
\[
\Theta(P;\mathcal M)=\{\theta(m):m\in\mathcal M,\ P_m=P\}.
\]
Point identification holds when this set is a singleton, or equivalently when $P_{m_1}=P_{m_2}$ implies $\theta(m_1)=\theta(m_2)$ for all compatible models. Sharp bounds mean bounds attained, or approached when the boundary is not attained, by compatible models. Writing this set is a definition, not a solution: the substantive task is to establish informative restrictions and derive or estimate the set. An unrestricted-confounding model may give only trivial bounds.

\textbf{Inference.} Identification, estimability and valid uncertainty quantification are separate questions. Fix $\alpha\in(0,1)$. A confidence procedure $C_n$ over a declared law class $\mathcal P$ has asymptotic uniform coverage at level $1-\alpha$ when
\[
\liminf_{n\to\infty}\inf_{P\in\mathcal P}\Pr_P\{\theta(P)\in C_n\}\geq1-\alpha.
\]
For set-identified targets the coverage event must be defined separately, for example coverage of the full identified set. If an entry uses identification-indexed models instead of uniquely indexed laws, the same coverage convention is applied over those models. Pointwise limits do not establish uniform validity, and asymptotic validity does not guarantee finite-sample precision. Sample sizes, dependence structures and weak-identification sequences are part of each application.

\textbf{Counterfactuals.} $Y_i(a)$ is a potential outcome under a specified intervention, including its timing, implementation and versions. It is not $Y_i$ conditional on voluntarily choosing $a$. A full intervention vector permits interference; shorthand scalar treatment notation does not silently impose no interference. Assignment, selection, attrition, anticipation and measurement must be specified. A claim about an instrument's local effect is not automatically a population policy effect.

\textbf{Economic equilibrium.} A counterfactual model includes best responses, constraints, feasibility, market clearing, state transitions and expectations consistent with the stated information. When several equilibria exist, either a declared selector is an input or the target is set-valued. Comparisons across policy regimes must use the stated selection convention. Reduced-form relationships are not presumed invariant after an equilibrium-changing intervention.

\textbf{Optimization and welfare.} Optimization is over the stated feasible set. If attainment is not established, $\sup$ or $\inf$ and $\varepsilon$-optimal choices replace an asserted maximizer or minimizer. For example, with value $V=\sup_{a\in\mathcal A}W(a)<\infty$, an $\varepsilon$-optimal set is $\{a\in\mathcal A:W(a)\geq V-\varepsilon\}$ for $\varepsilon>0$. Benefits, costs, risk and health or environmental outcomes can be aggregated only using declared units and conversion weights. Interpersonal utility weights, the treatment of fiscal transfers and foreign profits, discounting, and the behavioral welfare criterion are explicit inputs. A comparative-static derivative additionally requires existence and appropriate differentiability of the selected counterfactual.

\textbf{Sources.} Local labels [S1], [S2], and so forth restart in every question. Locators identify the relevant abstract, model, section or result at the level actually checked; a bibliographic or abstract-level verification is not represented as inspection of an entire proof. Working papers are distinguished from later publications and changing versions. Source summaries are paraphrased. All URL checks and editorial formulations refer to the audit date stated above.
% END ECONOMICS STATEMENT
\end{document}
